
\subsubsection{5.1 H-E1: Spurious Feature Dominance Exists}

\textbf{Finding:} Spurious features dominate from epoch 0, with attribution ratio 1.35---spurious regions receive 35\% more GradCAM activation than core regions from the first epoch.

\begin{table}[htbp]
\centering
\resizebox{\columnwidth}{!}{%
\begin{tabular}{lll}
\toprule
Epoch & Attribution Ratio (R) & Status \\
\midrule
0 & 1.348 & > 1.0 ✓ \\
5 & 1.273 & > 1.0 ✓ \\
10 & 1.431 & > 1.0 ✓ \\
13 (peak) & 1.590 & > 1.0 ✓ \\
25 & 1.255 & > 1.0 ✓ \\
50 & 1.255 & > 1.0 ✓ \\
\bottomrule
\end{tabular}
}
\end{table}

\textbf{Gate evaluation:} Dominance epoch < 10 required. Observed: dominance from epoch 0. \textbf{PASS.}

Spurious dominance is immediate (from initialization) and persistent (ratio never drops below 1.0). The ratio peaks at epoch 13 (1.59), suggesting maximum spurious reliance occurs early-to-mid training before stabilizing.

\textbf{Interpretation:} This confirms the simplicity bias phenomenon on Waterbirds. ERM-trained ResNet-50 exhibits spurious feature preference from the first forward pass, likely due to pretrained ImageNet features that respond more strongly to background textures than bird morphology.

\subsubsection{5.2 H-M1: Gradient Competition Hypothesis Falsified}

\textbf{Finding:} Contrary to the gradient competition hypothesis, spurious-aligned samples produce lower gradient norms than minority samples. The gradient ratio is 0.18, not >1.5 as predicted---an inversion by a factor of approximately 8.

\begin{table}[htbp]
\centering
\resizebox{\columnwidth}{!}{%
\begin{tabular}{lll}
\toprule
Seed & Early Epoch Ratio (1-10) & Direction \\
\midrule
42 & 0.175 & Minority > Spurious \\
123 & 0.173 & Minority > Spurious \\
456 & 0.181 & Minority > Spurious \\
**Mean** & **0.176 $\pm$ 0.004** & **Inverted** \\
\bottomrule
\end{tabular}
}
\end{table}

\textbf{Gate evaluation:} Ratio > 1.5 required. Observed: 0.18. \textbf{FAIL.}

The ratio starts around 0.45 at epoch 1 and decreases to 0.13-0.16 by epoch 50. The ratio is never above 1.0---minority samples produce stronger gradients at every epoch across all seeds.

Raw gradient norm data shows minority samples consistently produce 5-$7\times$ higher gradient norms than spurious-aligned samples throughout training. At epoch 1 (seed 42), spurious-aligned samples produce gradient norms of approximately 2.07e-4 while minority samples produce 4.60e-4. By epoch 50, these values decrease to 1.85e-5 and 1.60e-4 respectively, with the gap widening.

\textbf{Interpretation:} The gradient competition hypothesis is empirically falsified. Spurious features dominate (H-E1) but not because they receive stronger gradient signals. Spurious dominance arises from convergence speed---spurious patterns occupy simpler loss landscape regions that converge first, achieving low loss (and thus low gradients) quickly.

\subsubsection{5.3 Summary}

\begin{table}[htbp]
\centering
\resizebox{\columnwidth}{!}{%
\begin{tabular}{llll}
\toprule
Hypothesis & Prediction & Observed & Gate \\
\midrule
H-E1 (Existence) & R > 1.0 before epoch 10 & R = 1.35 at epoch 0 & **PASS** \\
H-M1 (Mechanism) & $\rho$ > 1.5 in epochs 1-10 & $\rho$ = 0.18 & **FAIL** \\
\bottomrule
\end{tabular}
}
\end{table}

Spurious dominance exists but the hypothesized mechanism is incorrect. The gradient competition hypothesis is falsified with high confidence (approximately $8\times$ inversion, consistent across 3 seeds).
